\SixthArticle{57}
Durch Kombination der Gleichungen 106, 107 erhalten wir demnach
\[
\begin{aligned}
{}[108]\quad&2AF(\alpha,\beta,\tfrac12,x)\\
&\quad=F(2\alpha,2\beta,\alpha+\beta+\tfrac12,\tfrac{1-\BareRoot{x}}2)\\&\qquad+F(2\alpha,2\beta,\alpha+\beta+\tfrac12,\tfrac{1+\BareRoot{x}}2)
\end{aligned}
\]
\[
\begin{aligned}
{}[109]\quad&2B\BareRoot{x}\mathbin{.}F(\alpha+\tfrac12,\beta+\tfrac12,\tfrac32,x)\\
&\quad=F(2\alpha,2\beta,\alpha+\beta+\tfrac12,\tfrac{1-\BareRoot{x}}2)\\&\qquad-F(2\alpha,2\beta,\alpha+\beta+\tfrac12,\tfrac{1+\BareRoot{x}}2)
\end{aligned}
\]
Ersetzen Sie in Gleichung 109 $\alpha$ durch $\alpha-\frac12$ und $\beta$ durch $\beta-\frac12$, so werden Sie leicht erkennen, dass sich ergibt:
\[
\begin{aligned}
{}[110]\quad&\frac{\alpha-\tfrac12\mathbin{.}\beta-\tfrac12}{\alpha+\beta-\tfrac12}A\BareRoot{x}\mathbin{.}F(\alpha,\beta,\tfrac32,x)\\
&\quad=F(2\alpha-1,2\beta-1,\alpha+\beta-\tfrac12,\tfrac{1+\BareRoot{x}}2)\\&\qquad-F(2\alpha-1,2\beta-1,\alpha+\beta-\tfrac12,\tfrac{1-\BareRoot{x}}2)
\end{aligned}
\]
